% GENERATED FILE. Edit canonical lessons, then run npm run generate:books.
% canonical-source-hash: fnv1a64:9d6471fc3eda9d08
% canonical-lessons: GU-C125-bhamar, GU-C125-dadhi, GU-C125-much, GU-C125-gardan, GU-C125-chati

\chapter{Eyebrow, Beard, Moustache, Neck, Chest}
\label{ch:gu-eyebrow-beard-moustache-neck-chest}
\hlchaptermodality{125}

\begin{chapteropening}
\textbf{By the end of this chapter:} \emph{I can say an eyebrow, a beard, a moustache, a neck and a chest.}

Complete the last lesson of chapter 125: I can say an eyebrow, a beard, a moustache, a neck and a chest.
\end{chapteropening}

\section[\texorpdfstring{\gu{ભમર} (bhamar)}{ભમર (bhamar)}]{\gu{ભમર} (bhamar) — an eyebrow}

\label{lesson:GU-C125-bhamar}



\begin{quote}
Before the new one: say the Gujarati for a cheek, then the Gujarati for a forehead.
\end{quote}

\subsection*{You'll want to know: \gu{ભમર}}
\textbf{\gu{ભમર}} (\emph{bhamar}) — \textquotedblleft{}an eyebrow\textquotedblright{}.

\begin{cousinweb}
\textbf{\gu{ભમર}} (\emph{bhamar}) is a tadbhava. CDIAL lists it, as \emph{bhammar}, under Sanskrit \emph{bhrū}, \textquotedblleft{}eyebrow\textquotedblright{} (CDIAL 9688). Its \emph{-m-} is not fully explained; CDIAL thinks it scarcely comes from the verb \emph{bhram}, \textquotedblleft{}to whirl\textquotedblright{}.
\end{cousinweb}

\begin{grammarlens}[title={What you've built}]
One more word for this chapter.
\end{grammarlens}

\subsection*{Your turn}
\begin{itemize}
\raggedright
  \item \emph{Say it:} \emph{bhamar}
  \item \emph{Say it:} \emph{bhamar}, once more
  \item \emph{Say it:} say \emph{bhamar}
  \item \emph{Recall:} say the Gujarati for a cheek, then the Gujarati for a forehead, then say \emph{bhamar} again
\end{itemize}

\begin{tcolorbox}[breakable,colback=teal!4,colframe=teal!35!black,title={Before you move on}]
What does \gu{ભમર} mean? (\textquotedblleft{}An eyebrow\textquotedblright{}.) Say it once more.
\end{tcolorbox}
\section[\texorpdfstring{\gu{દાઢી} (dāḍhī)}{દાઢી (dāḍhī)}]{\gu{દાઢી} (dāḍhī) — a beard, a chin}

\label{lesson:GU-C125-dadhi}



\begin{quote}
Before the new one: say the Gujarati for a forehead, then the Gujarati for an eyebrow.
\end{quote}

\subsection*{You'll want to know: \gu{દાઢી}}
\textbf{\gu{દાઢી}} (\emph{dāḍhī}) — \textquotedblleft{}a beard, a chin\textquotedblright{}.

\begin{cousinweb}
CDIAL lists \textbf{\gu{દાઢી}} (\emph{dāḍhī}), \textquotedblleft{}beard, chin\textquotedblright{}, under a reconstructed \emph{*dāṁṣṭra}, \textquotedblleft{}fang, tusk, beard\textquotedblright{}, beside Sanskrit \emph{daṁṣṭra}, \textquotedblleft{}tusk, fang\textquotedblright{} (CDIAL 6250). It holds that the sense could move from \textquotedblleft{}fang\textquotedblright{} to \textquotedblleft{}beard\textquotedblright{}, and does not rule out some Dravidian or Munda influence.
\end{cousinweb}

\begin{grammarlens}[title={What you've built}]
One more word for this chapter.
\end{grammarlens}

\subsection*{Your turn}
\begin{itemize}
\raggedright
  \item \emph{Say it:} \emph{dāḍhī}
  \item \emph{Say it:} \emph{dāḍhī}, once more
  \item \emph{Say it:} say \emph{dāḍhī}
  \item \emph{Recall:} say the Gujarati for a forehead, then the Gujarati for an eyebrow, then say \emph{dāḍhī} again
\end{itemize}

\begin{tcolorbox}[breakable,colback=teal!4,colframe=teal!35!black,title={Before you move on}]
What does \gu{દાઢી} mean? (\textquotedblleft{}A beard, a chin\textquotedblright{}.) Say it once more.
\end{tcolorbox}
\section[\texorpdfstring{\gu{મૂછ} (mūch)}{મૂછ (mūch)}]{\gu{મૂછ} (mūch) — a moustache}

\label{lesson:GU-C125-much}



\begin{quote}
Before the new one: say the Gujarati for an eyebrow, then the Gujarati for a beard.
\end{quote}

\subsection*{You'll want to know: \gu{મૂછ}}
\textbf{\gu{મૂછ}} (\emph{mūch}) — \textquotedblleft{}a moustache\textquotedblright{}.

\begin{cousinweb}
\textbf{\gu{મૂછ}} (\emph{mūch}) comes from \emph{*mucchā}, \textquotedblleft{}mustache\textquotedblright{}, a form CDIAL reconstructs (CDIAL 10180). CDIAL asks whether it blends \emph{gocchā}, \textquotedblleft{}furrow of the upper lip\textquotedblright{}, with \emph{mukha}, \textquotedblleft{}mouth\textquotedblright{}, or comes through Middle Indo-Aryan from Sanskrit \emph{śmaśru}, \textquotedblleft{}beard\textquotedblright{}.
\end{cousinweb}

\begin{grammarlens}[title={What you've built}]
One more word for this chapter.
\end{grammarlens}

\subsection*{Your turn}
\begin{itemize}
\raggedright
  \item \emph{Say it:} \emph{mūch}
  \item \emph{Say it:} \emph{mūch}, once more
  \item \emph{Say it:} say \emph{mūch}
  \item \emph{Recall:} say the Gujarati for an eyebrow, then the Gujarati for a beard, then say \emph{mūch} again
\end{itemize}

\begin{tcolorbox}[breakable,colback=teal!4,colframe=teal!35!black,title={Before you move on}]
What does \gu{મૂછ} mean? (\textquotedblleft{}A moustache\textquotedblright{}.) Say it once more.
\end{tcolorbox}
\section[\texorpdfstring{\gu{ગરદન} (gardan)}{ગરદન (gardan)}]{\gu{ગરદન} (gardan) — a neck}

\label{lesson:GU-C125-gardan}



\begin{quote}
Before the new one: say the Gujarati for a beard, then the Gujarati for a moustache.
\end{quote}

\subsection*{You'll want to know: \gu{ગરદન}}
\textbf{\gu{ગરદન}} (\emph{gardan}) — \textquotedblleft{}a neck\textquotedblright{}.

\begin{cousinweb}
\textbf{\gu{ગરદન}} (\emph{gardan}) is a loan, usually traced to Persian \emph{gardan}, \textquotedblleft{}neck\textquotedblright{}.
\end{cousinweb}

\begin{grammarlens}[title={What you've built}]
One more word for this chapter.
\end{grammarlens}

\subsection*{Your turn}
\begin{itemize}
\raggedright
  \item \emph{Say it:} \emph{gardan}
  \item \emph{Say it:} \emph{gardan}, once more
  \item \emph{Say it:} say \emph{gardan}
  \item \emph{Recall:} say the Gujarati for a beard, then the Gujarati for a moustache, then say \emph{gardan} again
\end{itemize}

\begin{tcolorbox}[breakable,colback=teal!4,colframe=teal!35!black,title={Before you move on}]
What does \gu{ગરદન} mean? (\textquotedblleft{}A neck\textquotedblright{}.) Say it once more.
\end{tcolorbox}
\section[\texorpdfstring{\gu{છાતી} (chātī)}{છાતી (chātī)}]{\gu{છાતી} (chātī) — a chest}

\label{lesson:GU-C125-chati}



\begin{quote}
Before the new one: say the Gujarati for a moustache, then the Gujarati for a neck.
\end{quote}

\subsection*{You'll want to know: \gu{છાતી}}
\textbf{\gu{છાતી}} (\emph{chātī}) — \textquotedblleft{}a chest\textquotedblright{}.

\begin{cousinweb}
\textbf{\gu{છાતી}} (\emph{chātī}) is a tadbhava, a word inherited and reshaped by sound change. It comes from \emph{*chātti}, \textquotedblleft{}chest, breast\textquotedblright{}, a form CDIAL reconstructs from its descendants (CDIAL 5014).
\end{cousinweb}

\begin{grammarlens}[title={What you've built}]
One more word for this chapter.
\end{grammarlens}

\subsection*{Your turn}
\begin{itemize}
\raggedright
  \item \emph{Say it:} \emph{chātī}
  \item \emph{Say it:} \emph{chātī}, once more
  \item \emph{Say it:} say \emph{chātī}
  \item \emph{Recall:} say the Gujarati for a moustache, then the Gujarati for a neck, then say \emph{chātī} again
\end{itemize}

\begin{tcolorbox}[breakable,colback=teal!4,colframe=teal!35!black,title={Before you move on}]
What does \gu{છાતી} mean? (\textquotedblleft{}A chest\textquotedblright{}.) Say it once more.
\end{tcolorbox}
